\begin{proof}
Sea $(x_i)_{1\leqslant i\leqslant m}$ un sistema de generadores de la
$A$-álgebra $C$; mostremos primero que se puede razonar por inducción sobre
$m$. Sea $C'$ la subálgebra de $C$ generada por
$x_1,\ldots,x_{m-1}$, y sea $\mathfrak n'=\mathfrak n\cap C'$.
El homomorfismo $A\to C_{\mathfrak n}$ se factoriza como
$A\to C'_{\mathfrak n'}\to C_{\mathfrak n}$, y está claro que
$A\to C'_{\mathfrak n'}$ y $C'_{\mathfrak n'}\to C_{\mathfrak n}$ son
homomorfismos locales; si $k''$ es el cuerpo residual de
$C'_{\mathfrak n'}$, $k\to k'$ se factoriza del mismo modo como
$k\to k''\to k'$, de modo que $k''$ es una extensión finita de $k$ y $k'$
una extensión finita de $k''$. La hipótesis de inducción implica que
$C'_{\mathfrak n'}$ es un $\boldsymbol P$-anillo; además, si
$S'=C'\setmin\mathfrak n'$, $C_{\mathfrak n}$ es un anillo local de
$S'^{-1}C$; como $C=C'[x_m]$, se tiene
$S'^{-1}C=C'_{\mathfrak n'}[x_m/1]$, y la hipótesis de inducción muestra
también que $C_{\mathfrak n}$ es un $\boldsymbol P$-anillo. La cuestión se reduce así al caso en que $C$ es una $A$-álgebra generada por
\emph{un solo} elemento $t$.

Apliquemos (7.4.3, (ii)); poniendo $A'=\widehat A$,
$B'=B\otimes_A A'$ es un anillo local de $C\otimes_A A'$ en un ideal
primo situado sobre $\mathfrak n$; como $A$ y $A'$ tienen el mismo cuerpo
residual $k$, el cuerpo residual de $C\otimes_A A'$ en ese ideal primo es
igual a $k'$. Por otra parte, $C\otimes_A A'$ es una $A'$-álgebra generada
por un solo elemento. Resulta, pues, de (7.4.3, (ii)) que es posible
limitarse a demostrar (7.4.4.1) cuando $\boldsymbol P$ es la propiedad
(iv$'$) de (7.3.8), $A$ es \emph{completo} y $C$ está generada por un solo
elemento $t$.

Para demostrar entonces que las fibras formales de $B=C_{\mathfrak n}$
son geométricamente regulares, apliquemos el criterio (7.3.16, \emph{b)}). Sea
$B_1$ una $B$-álgebra finita e \emph{íntegra}, generada por tanto por un
número finito de elementos enteros sobre $B$. Multiplicando estos
elementos por un elemento de $S=C\setmin\mathfrak n$, se puede suponer que son
enteros sobre $C$ y, por consiguiente, escribir $B_1=S^{-1}C_1$, donde
$C_1$ es una sub-$C$-álgebra de $B_1$ generada por un número finito de
elementos enteros sobre $C$; es, por tanto, una $C$-álgebra
\emph{finita e íntegra}. Por otra parte, $B_1$ es un anillo semilocal, y
todo anillo local $B_2$ de $B_1$ en un ideal maximal es un anillo local
de $C_1$ en un ideal primo, tal que $A\to B_2$ es un homomorfismo local;
además, el cuerpo residual de $B_2$ es una extensión finita de $k'$ y,
por tanto, de $k$. Se ve así, teniendo en cuenta (7.3.14) y (i), que
el problema se reduce a la cuestión siguiente: sea $C$ un anillo noetheriano
\emph{íntegro}, que contiene un subanillo $C_0$ que es una $A$-álgebra
generada por un solo elemento $t$, y tal que $C$ es una $C_0$-álgebra
\emph{finita}; si $\mathfrak n$ es un ideal maximal de $C$

\oldpage[IV-2]{201}

situado sobre el ideal maximal de $A$, y $B=C_{\mathfrak n}$, se trata
de demostrar que, para todo ideal primo $\mathfrak q$ de $\widehat B$
cuya imagen inversa en $B$ sea 0, el anillo
$\widehat B_{\mathfrak q}$ es \emph{regular}. Por otra parte, se puede
sustituir $A$ por su imagen en $C$, que es un anillo local completo
(como cociente de $A$) e íntegro. La conclusión que hay que demostrar
es entonces consecuencia del lema siguiente, en apariencia más general:
\end{proof}
